\documentclass[12pt]{exam}

\usepackage{mystyle}
\usepackage{listings}
\usepackage{xcolor}

\definecolor{codegreen}{rgb}{0,0.6,0}
\definecolor{codegray}{rgb}{0.5,0.5,0.5}
\definecolor{codepurple}{rgb}{0.58,0,0.82}
\definecolor{backcolour}{rgb}{0.95,0.95,0.92}

\lstdefinestyle{mystyle}{
    backgroundcolor=\color{backcolour},
    commentstyle=\color{codegreen},
    keywordstyle=\color{blue},
    numberstyle=\color{codegray},
    stringstyle=\color{codepurple},
    basicstyle=\ttfamily\footnotesize,
    breakatwhitespace=false,
    breaklines=true,
    captionpos=b,
    keepspaces=true,
    numbers=left,
    numbersep=5pt,
    showspaces=false,
    showstringspaces=false,
    showtabs=false,
    tabsize=2
}
\lstset{style=mystyle}

\title{MATH 3191 Exam 2 Review}

\date{}

\begin{document}
\maketitle
\section{Format of Exam}
The exam will take the following format
\begin{enumerate}
    \item Some short response questions that require a sentence or two explaining why a statement
        is true or false (A beginning of what a proof might look like)
    \item Longer response questions that will require you to compute a numerical answer and show work
        on how it was achieved
    \item Questions on lecture based content (No video exclusive content, but from the pdfs themselves)
\end{enumerate}
Note: Since this is a takehome exam, this exam will be fairly long (4-6 questions in each section. I haven't made it yet, so this is a rough estimate)
I will time this as if you have 3 hours to take it, but you have the entire week, so use your time wisely!
\section{Written Homework Problems}
You should be able to do all of the Written Homework problems from HW7 up to and including HW11. If you want
help with your solutions or how to approach these problems, just let me know!
\section{MyOpenMath Problems}
I would pay attention to the following topics associated with each problem below
\begin{enumerate}
    \item HW7: 1, 2, 4, 9, 10, 13, 19
        \begin{enumerate}
            \item Be able to describe requirements for a vector space and demonstrate its properties
            \item Determine if a given set is a subspace of a known vector space ($\R^n,\mathcal{P}(\R)$)
            \item Describe what a null space/kernel is
            \item Describe what a column space is
            \item Describe what a spanning list is
        \end{enumerate}
    \item HW8: 1, 2, 3, 6, 7, 8, 10, 13, 18
        \begin{enumerate}
            \item Know the requirements for a basis of a vector space
            \item Determine a basis for a null space and column space
            \item Determine the $\mathcal{B}$-coordinates of a given vector for a particular basis
            \item Given the $\mathcal{B}$-coordinates of a vector in a particular basis, determine the vector in the standard basis
            \item Know how the dimension of the null space and column space relate.
        \end{enumerate}
    \item HW9: 1, 3, 4, 6, 7, 10, 11, 14, 16, 18
        \begin{enumerate}
            \item Determine the change of basis matrix for two given bases
            \item Convert a coordinate vector from one basis to another
            \item Determine if a given vector is an eigenvector
            \item Determine the eigenvalues of a matrix
            \item Determine the eigenvalue of a matrix power
        \end{enumerate}
    \item HW10: 1, 2, 5, 6, 9, 11, 15, 18, 21
        \begin{enumerate}
            \item Diagonalize a matrix with real eigenvalues
            \item Compute a large power of a matrix ($A^n$ for $n\geq 4$)
            \item Know the requirements for a matrix to be diagonalizable and properties of diagonal matrices
            \item Compute complex eigenvalues of a $2\times 2$ or $3\times 3$ matrix
            \item These last objectives are only if we get to them on 4-03 in class
                \begin{enumerate}
                    \item Determine if a given matrix is a transition matrix
                    \item Compute the long run probabilities of Markov Chain
                \end{enumerate}
        \end{enumerate}
    \item HW11:  3, 4, 5, 6, 7, 8, 9
        \begin{enumerate}
            \item Compute inner products of vectors
            \item Determine how to make a given vector orthogonal to another
            \item Determine Orthogonal Complements of a set
            \item Determine if a given set of vectors are orthogonal, orthonormal, or neither
            \item Properties of Orthonormal Lists
        \end{enumerate}
\end{enumerate}
\section{Lecture Based Problems}
\begin{enumerate}
    \item Explain what a similarity transformation is and what it does. IE: given $A = CBC^{-1}$, describe
        what each step of $CBC^{-1}\vec{x}$ is doing.
    \item Explain why using a different basis can make some problems easier (Like computing $A^n$).
    \item Explain what information diagonalization gives us about a matrix
    \item Be able to explain what a Markov chain is intuitively and algebraically
    \item Explain what an Inner product is, and be able to demonstrate a simple example is an IP
    \item Use a given inner product to test orthogonality
    \item Explain what Orthogonality and Orthonormality are along with theorems in slides.
\end{enumerate}
\section{Python Lab Problems}
You should be able to explain what every function we have been taught in the labs do, and translate python
code into a plain english explanation of what it does. For example consider the following code snippet
\begin{lstlisting}[language=Python]
import numpy as np
A = np.array([[1, 2, 4],
              [2, 1, 7],
              [1, 1, 1]])
b = np.array([1, 1, 1])
x = np.linalg.solve(A,b)
\end{lstlisting}
I would expect a response like ``This code solves the system
    \[
        \begin{bmatrix}1&2&4 \\ 2&1&7 \\ 1&1&1 \end{bmatrix}\begin{bmatrix}x_1\\x_2\\x_3\end{bmatrix} =
            \begin{bmatrix} 1\\1\\1\end{bmatrix}
    \]
        for $\vec{x} = \begin{bmatrix}x_1\\x_2\\x_3\end{bmatrix}$ using numpy''.
    You would NOT have to then solve this system unless explicitly asked to.

    You should also be able to tell me what kind of output to expect. For example in the above snippet,
    I would expect an answer similar to ``you will get an array of 3 numbers if a solution exists''

    I will NOT ask you to write code by hand, only to do the above tasks.
\end{document}
